\documentclass[12pt]{article}
\usepackage[utf8]{inputenc}
\usepackage{amsmath, amssymb, amsthm}
\usepackage{geometry}

% Set margins to match a standard homework document
\geometry{a4paper, margin=1in}

\title{Challenge Problem Set 6}
\parskip=1ex
\author{Math 292 Spring 2026 \\ Rutgers University}
\date{}

\begin{document}

\maketitle

\noindent{}\textbf{Exercise 1:}

Suppose $A$ is a self-adjoint matrix with non-positive eigenvalues.

Then, by the Spectral Theorem, $A$ can be diagonalized into an orthonormal basis. That is
\begin{align*}
    A = UDU^*,
\end{align*}
where $U$ is a unitary matrix and $D$ is a diagonal matrix.


Thus, we also have
\begin{align*}
    e^{tA} = Ue^{tD}U^*.
\end{align*}

For any $\mathbf{x}$, we can rewrite it as a linear combination of this orthonormal basis.
That is, for constants $\mathbf{c}$, we have
\begin{align*}
    ||\mathbf{x}||^2 & = ||U\mathbf{c}||^2          \\
                     & = <U\mathbf{c}, U\mathbf{c}> \\
                     & = ||\mathbf{c}||^2
\end{align*}
Let $A$ have $n$ columns. Then
\begin{align*}
    ||\mathbf{x}||^2 = \sum_{i=1}^{n} c_i^2.
\end{align*}

Now we consider
\begin{align*}
    ||e^{tA}\mathbf{x}||^2 & = ||Ue^{tD}U^*\mathbf{x}||^2           \\
                           & = ||Ue^{tD}\mathbf{c}||^2              \\
                           & = <e^{tD}\mathbf{c}, e^{tD}\mathbf{c}> \\
                           & = \sum_{i=1}^n e^{2t\lambda_i}c_i^2.
\end{align*}
Since $t \geq 0$ and $\lambda_i \leq 0$, it follows that $t\lambda_i \leq 0$ and $e^{2t\lambda_i} \leq e^{0} = 1$.
Thus,
\begin{align*}
    ||e^{tA}\mathbf{x}||^2 \leq \sum_{i=1}^n c_i^2 = ||\mathbf{x}||^2.
\end{align*}
Taking the square root of both sides gives
\begin{align*}
    ||e^{tA}\mathbf{x}|| \leq ||\mathbf{x}||.
\end{align*}

Therefore, we conclude that $e^{tA}$ is contractive if $A$ has all non-positive eigenvalues.

\noindent{}\textbf{Exercise 2:}

Suppose $K$ is a skew-adjoint square matrix.
Then, we compute
\begin{align*}
    (iK)^* = -i(-K) = iK.
\end{align*}
That is, $iK$ is a self-adjoint matrix.

Thus, by Spectral Theorem, the eigenvalues of $K$ must be purely imaginary and
$K$ can be diagonalized into an orthonormal basis $UDU^*$. Thus, $e^{tK} = Ue^{tD}U^*$.
Consider the adjoint of $e^{tK}$:
\begin{align*}
    (e^{tK})^* & = (Ue^{tD}U^*)^* \\
               & = Ue^{-tD}U^*    \\
               & = e^{-tK}.
\end{align*}
We can check that this is exactly the inverse:
\begin{align*}
    e^{tK} (e^{tK})^* = e^{tK} e^{-tK} = I = e^{tK} (e^{tK})^{-1}.
\end{align*}
Thus, we have that $e^{tK}$ is unitary.

To show $e^{tK}$ is also norm-preserving, we compute
\begin{align*}
    ||e^{tK}\mathbf{x}||^2 & = <e^{tK}\mathbf{x}, e^{tK}\mathbf{x}>      \\
                           & = \mathbf{x}^T (e^{tK})^* e^{tK} \mathbf{x} \\
                           & = \mathbf{x}^T I \mathbf{x}                 \\
                           & = \mathbf{x}^T \mathbf{x}                   \\
                           & = ||\mathbf{x}||^2.
\end{align*}
Taking the square root of both sides gives $||e^{tK}\mathbf{x}|| = ||\mathbf{x}||$.
Hence $e^{tK}$ is norm-preserving.


\noindent{}\textbf{Exercise 3:}

Let $A = \frac{1}{2}(A + A^*) + \frac{1}{2}(A - A^*) =: H + K$.

Suppose $H$ has non-positive eigenvalues.

Taking the exponential of $A=H+K$, we have
\begin{align*}
    e^{tA} & = e^{t(H+K)}.
\end{align*}

Then, applying the trotter product, we get
\begin{align*}
    e^{t(H+K)} = \lim_{n\rightarrow{\infty}}(e^{\frac{t}{n}H}e^{\frac{t}{n}K})^n.
\end{align*}

Since $T_n \rightarrow T \implies ||T_n|| \rightarrow ||T||$ under the operator norm, we have
\begin{align*}
    ||e^{t(H+K)}|| & = ||\lim_{n\rightarrow{\infty}}(e^{\frac{t}{n}H}e^{\frac{t}{n}K})^n||  \\
                   & = \lim_{n\rightarrow{\infty}} ||(e^{\frac{t}{n}H}e^{\frac{t}{n}K})^n||
\end{align*}

Consider $(e^{\frac{t}{n}H}e^{\frac{t}{n}K})^n$, because $||AB|| \leq ||A||~||B||$ under the operator norm, we have
\begin{align*}
    ||(e^{\frac{t}{n}H}e^{\frac{t}{n}K})^n|| & \leq ||(e^{\frac{t}{n}H}e^{\frac{t}{n}K})||^n        \\
                                             & \leq ||e^{\frac{t}{n}H}||^n ~ ||e^{\frac{t}{n}K}||^n
\end{align*}

Since $H$ is also self-adjoint, it follows that $||e^{tH}|| \leq 1$.
Since $K$ is skew-adjoint, it follows that $||e^{tK}|| = 1$.

Thus,
\begin{align*}
    ||e^{\frac{t}{n}H}||^n ~ ||e^{\frac{t}{n}K}||^n \leq 1^n \cdot 1^n = 1.
\end{align*}

Hence,
\begin{align*}
    ||e^{tA}|| = \lim_{n\rightarrow{\infty}} ||(e^{\frac{t}{n}H}e^{\frac{t}{n}K})^n|| \leq \lim_{n\rightarrow{\infty}} 1 = 1.
\end{align*}
That is, $||e^{tA}||  \leq 1$.

Therefore, we conclude that $e^{tA}$ is contractive if $H$ has only non-positive real eigenvalues.

\noindent{}\textbf{Exercise 4:}

Suppose that $A$ is any constant square matrix and that $e^{tA}$ is contractive for all $t \geq 0$.
By definition of contractive, we have
\begin{align*}
    ||e^{tA}x||^2 \leq ||x||^2 \\
    ||e^{tA}x||^2 - ||x||^2 \leq 0
\end{align*}
Notice that this is equal to $||e^{tA}x||^2 - ||e^{0}x||^2 \leq 0$ since $e^{0}=I$.
Then, dividing through by t and taking $t \rightarrow 0$, we have
\begin{align*}
    \lim_{t \rightarrow 0} \frac{||e^{tA}x||^2 - ||x||^2}{t} \leq 0.
\end{align*}
Notice this is exactly the definition of the derivative. So,
\begin{align*}
    \frac{d}{dt} ||e^{tA}x||^2 ~\big|_{t=0} \leq 0.
\end{align*}

Computing the derivative of $||e^{tA}x||$, we have
\begin{align*}
    \frac{d}{dt} ||e^{tA}x||^2 & = \langle Ae^{tA}x, e^{tA}x \rangle + \langle e^{tA}x, Ae^{tA}x \rangle.
\end{align*}
Evaluating this at $t=0$ gives
\begin{align*}
    \frac{d}{dt} ||e^{tA}x||^2 ~\big|_{t=0} & = \langle Ax, x \rangle + \langle x, Ax \rangle \\
                                            & = 2Re\langle x, Ax \rangle.
\end{align*}
Using the previous equation $\frac{d}{dt} ||e^{tA}x||^2 ~\big|_{t=0} \leq 0$, this reduces to
\begin{align*}
    Re(x, Ax) \leq 0
\end{align*}

Therefore, we conclude that $A$ is dissipative.

\noindent{}\textbf{Exercise 5:}

Consider the following matrices:
\begin{align*}
    A_1 = \begin{bmatrix}
              -\frac{1}{3} & 1 \\ 0 & -\frac{1}{3}
          \end{bmatrix}
    \hspace{3ex}
    A_2 = \begin{bmatrix}
              -1 & 1 \\ 0 & -1
          \end{bmatrix}
\end{align*}
Notice that $A_1$ has eigenvalues $\lambda_1 = \lambda_2 = -\frac{1}{3}$
and $A_2$ has eigenvalues $\mu_1 = \mu_2 = -1$.

Consider the associated self-adjoint part of $A_1$, that being $H_1$.
\begin{align*}
    H_1 & = \frac{1}{2}(A_1 + A_1^*)        \\
        & = \frac{1}{6} \begin{bmatrix}
                            -2 & 3 \\ 3 & -2
                        \end{bmatrix}
\end{align*}
Then, we solve for its determinant:
\begin{align*}
    det(H_1) & = -\frac{5}{36} < 0.
\end{align*}
Recall that $det(H_1) = \lambda_1 \cdot \lambda_2$.
Thus, $H_1$ has two eigenvalues of opposite signs.
From exercises 3 and 4, we concluded that $H$ has negative eigenvalues if and only if $A$ is dissipative.
Since $H_1$ has a positive eigenvalue, it follows that $A_1$ isn't dissipative.
It follows that $e^{tA_1}$ is not contractive as well.

Next, we consider the associated self-adjoint part of $A_2$, that being $H_2$.
\begin{align*}
    H_2 & = \frac{1}{2}(A_2 + A_2^*)                \\
        & = \begin{bmatrix}
                -1 & \frac{1}{2} \\ \frac{1}{2} & -1
            \end{bmatrix}.
\end{align*}
Then, solving for its determinant:
\begin{align*}
    det(H_2) & = \frac{3}{4} > 0.
\end{align*}
So $H_2$'s eigenvalues are of the same sign.
Notice that $tr(H_2) = -2$ is negative.
Hence, both eigenvalues must also be negative.
Since $H_2$ has only non-positive eigenvalues, we conclude that $A_2$ is dissipative, and so $e^{tA_2}$ is contractive.


\end{document}

